\section{Reversibilidad lógica y principio de Landauer en las fibras de TPK}
\label{sec:informacion-fibra-landauer-rev7}
\index{Landauer}\index{información!fibra}

Sea \((\Omega,2^\Omega,\mu)\) un espacio probabilizado finito, sea
\(X:\Omega\to\mathcal X\) una variable aleatoria con valores en el espacio
\(\mathcal X\) de estados completos y sea \(q:\mathcal X\to Y\) un lector.
La regla de la cadena da
\[
 \mathsf H(X)
 =
 \mathsf H(q(X))+\mathsf H(X\mid q(X)).
\]
El segundo sumando mide la información de la fibra que el lector no
publica.

\begin{theorem}[Landauer nulo para el transporte completo y coste de la proyección]
\label{thm:landauer-relativo-lector-rev7}
Si \(U:\mathcal X\to\mathcal X\) es una actualización biyectiva del estado
HMT completo,
entonces
\[
 \boxed{\mathsf H(X\mid U(X))=0.}
\]
Para una salida proyectada \(q(U(X))\), la información descartada es
\[
 \mathsf H(X\mid q(U(X))).
\]
Un reset físico que identifique estados distintos de la memoria debe
contabilizar esa información en el entorno.
\end{theorem}

\begin{proof}
La biyectividad implica
\(X=U^{-1}(U(X))\); por tanto \(X\) queda determinado por \(U(X)\) y la
entropía condicional es cero. Para el lector \(q\), la regla de la cadena
separa la entropía publicada de la entropía de la fibra. Si la operación de
reinicio no es inyectiva, la salida ya no permite reconstruir el estado de
memoria. El
principio termodinámico de Landauer acota el coste de la implementación en
función de esa entropía borrada; la actualización reversible no aporta por
sí sola un término de borrado.
\end{proof}

En nats y bits, respectivamente,
\[
 Q\ge k_B\Theta\,\mathsf H_{\rm nat},
 \qquad
 Q\ge k_B\Theta\log 2\,\mathsf H_{\rm bit}.
\]
Para el transporte biyectivo del estado completo, la entropía de borrado es
cero y, por tanto, el término mínimo de Landauer asociado exclusivamente al
borrado satisface
\[
\boxed{Q_{\rm L}^{\min}[U]=k_B\Theta\,
 \mathsf H(X\mid U(X))=0.}
\]
Para el reinicio de un TRIT uniforme, en cambio,
\[
 \boxed{Q_{\rm reset}\ge k_B\Theta\log3.}
\]
Estas dos ecuaciones localizan el coste en la pérdida de información: el
coste aparece cuando la proyección o el reinicio olvidan la fibra.
La preparación, la proyección y el reinicio pueden producir entropía; el
transporte reversible conserva la información genealógica. El valor nulo es
el límite lógico de Landauer para esa actualización completa, mientras los
costes dinámicos de control, velocidad finita y acoplamiento al entorno
pertenecen a la realización física del dispositivo.

\begin{statusbox}[title={Alcance del Landauer nulo}]
La igualdad \(\mathsf H(X\mid U(X))=0\) es exacta para una actualización
biyectiva del estado enriquecido. La cota
\(Q_{\rm reset}\ge k_B\Theta\log3\) es la especialización de Landauer al
borrado de una decisión triádica uniforme. La realización de ambas fórmulas
en joules utiliza el transductor de Boltzmann y una temperatura física
declarada.
\end{statusbox}
